\section{Evaluation}
We compare four proving runs, which we call \emph{cells} (Table~1). Cells~1--3 all prove PLUM at $\lambda=80$. Cell~1 runs the Griffin hash without the accelerator, Cell~2 with it, and Cell~3 is a SHA-3 control. We use $\lambda=80$ because it is the highest level that completes on this hardware. The full $128$ level is untested and expected to run out of memory. Cell~4 (Aurora~\cite{cryptoeprint:2018/828} proving Loquat~\cite{cryptoeprint:2024/868}) is a different scheme over different numbers and is the only one that hides its inputs, so it serves only as a rough cross-scheme reference.

{\par\medskip\centering\footnotesize
\textbf{Table 1.}~Four cells ($24$\,GB; cells 1--3: PLUM, $\lambda{=}80$, $n{=}5$ for cells 2--3; cell 4: Loquat in Aurora, target $80$-bit, $n{=}3$).\\[3pt]
\begin{tabular}{cll}
\hline
Cell & Setting & Prove \\
\hline
1 & no precompile      & $>5$\,h \\
2 & Griffin precompile & $14.24$\,min \\
3 & SHA-3 control      & $13.24$\,min \\
4 & Aurora (Loquat)    & ${\approx}22$\,min \\
\hline
\end{tabular}\par\medskip}

With the accelerator, verification finishes for the first time in the study (using at most ${\approx}13.7$\,GiB). Cell~1 stayed within memory but was stopped, unfinished, after five hours. On cost, the control is the sharpest result. Cell~3 is a diagnostic only, not a secure form of PLUM, because PLUM's security analysis does not cover SHA-3. Even with the accelerator, PLUM verification still proves $1.08\times$ slower, and runs about $1.11\times$ more steps, than the plain software SHA-3, so a standard hash beats the accelerated one. The criterion predicts this. The accelerator removes the emulation cost, but its own circuit must still be proved. A separate control, in which the same multiplication carried out in one part rather than seven costs $164\times$ less, confirms the penalty comes from splitting numbers across parts. Both halves of the credential, issuing it and showing it, also run and prove on the zkVM (${\approx}69$\,min and ${\approx}103$--$137$\,min), so the study covers the whole credential, not one signature alone.
